\section{Regime 1 --- Independent Agents (Multiplication Law)}
\label{sec:regime1}

Throughout this section, let $(A_1, A_2, \ldots, A_n)$ be a sequential pipeline of
agents as in Definition~\ref{def:pipeline}, with $A_i \sim \mathrm{Bernoulli}(p_i)$.
Unless stated otherwise, the agents are assumed \textbf{mutually independent}.

%% ============================================================
\subsection{The Multiplication Law}
%% ============================================================

\begin{theorem}[Multiplication Law for Independent Pipelines]
\label{thm:mult-law}
If agents $A_1, \ldots, A_n$ are mutually independent with success probabilities
$p_1, \ldots, p_n$, then
\[
  P_{\mathrm{sys}}(n) \;=\; \prod_{i=1}^{n} p_i.
\]
\end{theorem}

\begin{proof}
By Definition~\ref{def:psys},
\[
  P_{\mathrm{sys}}(n) = P\!\Bigl(\,\bigcap_{i=1}^{n} \{A_i = 1\}\Bigr).
\]
Since the $A_i$ are mutually independent, the probability of the intersection of
independent events equals the product of their probabilities (this is the defining
property of mutual independence applied to the events $\{A_i = 1\}$):
\[
  P\!\Bigl(\,\bigcap_{i=1}^{n} \{A_i = 1\}\Bigr)
  \;=\; \prod_{i=1}^{n} P(A_i = 1)
  \;=\; \prod_{i=1}^{n} p_i.
  \qedhere
\]
\end{proof}

%% ============================================================
\subsection{Exponential Decay of Reliability}
%% ============================================================

\begin{theorem}[Exponential Decay]
\label{thm:exp-decay}
Let all agents have uniform capability $p \in (0,1)$. Then
\[
  P_{\mathrm{sys}}(n) \;=\; p^n \;=\; e^{\,n \ln p},
\]
which decays exponentially in $n$ at rate $|\ln p| = -\ln p > 0$.

More precisely, for any target reliability $P_{\mathrm{sys}}(n) \geq \delta$ with
$\delta \in (0,1)$, the maximum number of agents is
\[
  n_{\max}(\delta, p) \;=\; \Bigl\lfloor \frac{\ln \delta}{\ln p} \Bigr\rfloor.
\]
\end{theorem}

\begin{proof}
By Theorem~\ref{thm:mult-law} with $p_i = p$ for all $i$,
\[
  P_{\mathrm{sys}}(n) = p^n.
\]
Since $0 < p < 1$, we have $\ln p < 0$, so $p^n = e^{n \ln p}$ is a strictly
decreasing exponential function of $n$. The decay rate per agent is
\[
  -\frac{d}{dn}\ln P_{\mathrm{sys}}(n) = -\ln p = |\ln p| > 0.
\]

For the bound, require $p^n \geq \delta$. Taking logarithms (note $\ln p < 0$, so
the inequality reverses):
\[
  n \ln p \geq \ln \delta
  \;\;\Longleftrightarrow\;\;
  n \leq \frac{\ln \delta}{\ln p}.
\]
Since $n$ must be a positive integer, $n_{\max} = \lfloor \ln\delta / \ln p \rfloor$.
\end{proof}

\begin{corollary}[Half-Life of a Pipeline]
\label{cor:half-life}
The number of agents at which system reliability drops to half the single-agent
capability is
\[
  n_{1/2} \;=\; \Bigl\lfloor \frac{\ln(1/2)}{\ln p} \Bigr\rfloor + 1
  \;=\; \Bigl\lfloor \frac{\ln 2}{|\ln p|} \Bigr\rfloor + 1.
\]
For high-capability agents ($p$ close to 1), using $\ln p \approx -(1-p)$:
\[
  n_{1/2} \;\approx\; \frac{\ln 2}{1-p}.
\]
\end{corollary}

\begin{proof}
We seek the smallest $n$ such that $p^n \leq p/2$ (half the single-agent reliability).
Dividing by $p$, this is $p^{n-1} \leq 1/2$, i.e., $(n-1) \geq \ln(1/2)/\ln p =
\ln 2 / |\ln p|$. The approximation follows from the first-order Taylor expansion
$\ln p = \ln(1-(1-p)) \approx -(1-p)$ for $p$ near 1.
\end{proof}

%% ============================================================
\subsection{Heterogeneous Agents: AM--GM Bound}
%% ============================================================

\begin{lemma}[Uniform Capability Maximizes System Reliability]
\label{lem:am-gm}
Among all pipelines of $n$ independent agents with a fixed sum of capabilities
$\sum_{i=1}^{n} p_i = S$ (where $S \leq n$), the system success probability
$P_{\mathrm{sys}} = \prod p_i$ is maximized when $p_i = S/n$ for all $i$.

Equivalently, for any fixed arithmetic mean $\bar{p} = S/n$:
\[
  \prod_{i=1}^{n} p_i \;\leq\; \bar{p}^{\,n},
\]
with equality iff $p_1 = p_2 = \cdots = p_n$.
\end{lemma}

\begin{proof}
This is the AM--GM inequality applied in logarithmic form. Since $\ln$ is concave,
Jensen's inequality gives
\[
  \frac{1}{n}\sum_{i=1}^{n} \ln p_i \;\leq\; \ln\!\Bigl(\frac{1}{n}\sum_{i=1}^{n} p_i\Bigr) = \ln \bar{p},
\]
hence $\sum \ln p_i \leq n \ln \bar{p}$, i.e., $\prod p_i \leq \bar{p}^n$.
Equality in Jensen's inequality holds iff all $p_i$ are equal.
\end{proof}

%% ============================================================
\subsection{Critical Threshold}
%% ============================================================

\begin{theorem}[Critical Threshold for Beneficial Agent Addition]
\label{thm:critical-threshold}
Consider a uniform-capability pipeline with per-link coordination cost $c \in [0,1)$.
The effective system success probability is (Definition~\ref{def:coord-cost}):
\[
  P_{\mathrm{eff}}(n) \;=\; p^n \cdot (1-c)^{n-1}.
\]

Adding an $(n+1)$-th agent is beneficial (i.e., $P_{\mathrm{eff}}(n+1) > P_{\mathrm{eff}}(n)$) if
and only if
\[
  p(1-c) > 1.
\]

Therefore the \emph{critical threshold} is
\[
  \boxed{p^* \;=\; \frac{1}{1-c}.}
\]

\begin{itemize}
  \item For $p > p^*$: every additional agent \emph{increases} effective system reliability.
  \item For $p < p^*$: every additional agent \emph{decreases} effective system reliability.
  \item For $p = p^*$: $P_{\mathrm{eff}}$ is constant in $n$.
\end{itemize}

Note: Since $p \leq 1$ and $p^* = 1/(1-c)$, the condition $p > p^*$ can only be
satisfied when $c = 0$ and $p = 1$. For any $c > 0$, we have $p^* > 1$, which means
$p < p^*$ always holds --- adding agents to a pipeline with coordination cost always
degrades performance. This is the fundamental bottleneck of sequential architectures.
\end{theorem}

\begin{proof}
Compute the ratio:
\[
  \frac{P_{\mathrm{eff}}(n+1)}{P_{\mathrm{eff}}(n)}
  \;=\; \frac{p^{n+1}(1-c)^{n}}{p^{n}(1-c)^{n-1}}
  \;=\; p(1-c).
\]
This ratio is independent of $n$. We have:
\begin{align*}
  P_{\mathrm{eff}}(n+1) > P_{\mathrm{eff}}(n)
  &\;\Longleftrightarrow\; p(1-c) > 1
  \;\Longleftrightarrow\; p > \frac{1}{1-c} = p^*.
\end{align*}

Since $c \in [0,1)$, we have $1-c \in (0,1]$, hence $p^* = 1/(1-c) \geq 1$.
Because $p \leq 1$ by definition, the condition $p > p^*$ requires $p^* < 1$,
which is impossible. When $c = 0$, $p^* = 1$ and the condition becomes $p > 1$,
never satisfied. Thus for $c > 0$, $p^* > 1 > p$ strictly, and adding agents
always reduces $P_{\mathrm{eff}}$.
\end{proof}

\begin{remark}[Interpretation for Zero Coordination Cost]
\label{rmk:zero-cost}
When $c = 0$ (no coordination overhead), the effective probability reduces to
$P_{\mathrm{eff}}(n) = p^n$, which is still decreasing for $p < 1$. The critical threshold
$p^* = 1$ means that only a perfect agent ($p = 1$) can be added without degradation.
This confirms that the sequential pipeline is inherently fragile: it is the
multiplicative structure itself, not just coordination cost, that drives exponential
decay.
\end{remark}

%% ============================================================
\subsection{Optimal Number of Agents}
%% ============================================================

The preceding theorem shows that in a pure sequential pipeline, adding agents always
hurts when $c > 0$. This means the trivially optimal pipeline length is $n^* = 1$.
However, this result assumes agents perform identical tasks. A more realistic model
assumes that the task \emph{requires} decomposition: a task of complexity $T$ must be
divided among agents, and more agents means each agent handles a simpler subtask,
increasing its per-subtask success probability.

\begin{definition}[Task Decomposition Model]
\label{def:task-decomp}
A task of complexity $T > 0$ is divided among $n$ agents. Each agent handles a subtask
of complexity $T/n$. The per-agent success probability is a decreasing function of
subtask complexity:
\[
  p(n) \;=\; g\!\left(\frac{T}{n}\right),
\]
where $g : [0, \infty) \to [0,1]$ is a decreasing function with $g(0) = 1$ and
$\lim_{x \to \infty} g(x) = 0$. A natural choice is the exponential model:
\[
  p(n) \;=\; e^{-T/n}.
\]
\end{definition}

\begin{theorem}[Optimal Pipeline Length --- Exponential Decomposition]
\label{thm:optimal-n}
Under the exponential task decomposition model $p(n) = e^{-T/n}$ with per-link
coordination cost $c \in (0,1)$, the effective system probability is
\[
  P_{\mathrm{eff}}(n) \;=\; e^{-T} \cdot (1-c)^{n-1}.
\]
This is maximized at $n^* = 1$ for all $T > 0$ and $c > 0$.
\end{theorem}

\begin{proof}
Substitute $p(n) = e^{-T/n}$ into $P_{\mathrm{eff}}$:
\[
  P_{\mathrm{eff}}(n)
  = \bigl(e^{-T/n}\bigr)^n \cdot (1-c)^{n-1}
  = e^{-T} \cdot (1-c)^{n-1}.
\]
The factor $e^{-T}$ is independent of $n$. Since $(1-c) < 1$, the factor $(1-c)^{n-1}$
is strictly decreasing in $n$. Hence $P_{\mathrm{eff}}$ is maximized at $n = 1$, where
$P_{\mathrm{eff}}(1) = e^{-T}$.
\end{proof}

\begin{remark}
Theorem~\ref{thm:optimal-n} reveals a degenerate case: the exponential decomposition
$p(n) = e^{-T/n}$ satisfies $[p(n)]^n = e^{-T}$ identically, so task decomposition
yields \emph{no net benefit} to the multiplicative product. All benefit from simpler
subtasks is exactly canceled by having more agents. Coordination cost then strictly
favors fewer agents. To obtain a non-trivial optimum we need decomposition that
yields a net benefit exceeding coordination cost.
\end{remark}

\begin{definition}[Power-Law Decomposition Model]
\label{def:power-decomp}
A more favorable decomposition model assumes per-agent success probability
\[
  p(n) \;=\; 1 - \left(\frac{T}{n}\right)^{\!\beta}
\]
for some $\beta > 0$ and $T/n < 1$ (i.e., $n > T$). The parameter $\beta$ controls
how rapidly success probability improves as subtask complexity decreases.
\end{definition}

\begin{theorem}[Optimal Pipeline Length --- General Decomposition]
\label{thm:optimal-n-general}
Let $p(n)$ be a decomposition model such that the log-effective probability
\[
  \ell(n) \;:=\; \ln P_{\mathrm{eff}}(n) \;=\; n \ln p(n) + (n-1)\ln(1-c)
\]
is a concave function of $n$ (treating $n$ as continuous). Then the optimal pipeline
length $n^*$ satisfies
\[
  \boxed{\ln p(n^*) + n^* \cdot \frac{p'(n^*)}{p(n^*)} + \ln(1-c) = 0.}
\]
Equivalently,
\[
  n^* \cdot \frac{p'(n^*)}{p(n^*)} \;=\; -\ln p(n^*) - \ln(1-c)
  \;=\; -\ln\!\bigl[p(n^*)(1-c)\bigr].
\]
\end{theorem}

\begin{proof}
Differentiate $\ell(n)$ with respect to $n$ (treating $n$ as continuous):
\[
  \ell'(n) = \ln p(n) + n \cdot \frac{p'(n)}{p(n)} + \ln(1-c).
\]
Setting $\ell'(n^*) = 0$ gives the stated equation. If $\ell$ is concave, this
critical point is a global maximum.
\end{proof}

\begin{corollary}[Explicit Solution for Near-Perfect Agents]
\label{cor:optimal-approx}
When $p$ is close to 1 (high capability), write $p(n) = 1 - \epsilon(n)$ where
$\epsilon(n) \ll 1$. Then $\ln p(n) \approx -\epsilon(n)$ and $p'(n)/p(n) \approx
-\epsilon'(n)$. The optimality condition becomes
\[
  -\epsilon(n^*) - n^* \epsilon'(n^*) + \ln(1-c) = 0.
\]
For the power-law model $\epsilon(n) = (T/n)^\beta$ with $\epsilon'(n) =
-\beta T^\beta / n^{\beta+1}$:
\[
  -\frac{T^\beta}{(n^*)^\beta} + \frac{\beta T^\beta}{(n^*)^\beta} + \ln(1-c) = 0,
\]
giving
\[
  \boxed{n^* \;=\; T \cdot \left(\frac{\beta - 1}{|\ln(1-c)|}\right)^{\!1/\beta}}
  \quad (\text{for } \beta > 1).
\]
For $\beta \leq 1$, no finite optimum exists (the numerator is non-positive), meaning
decomposition benefit is too weak to overcome coordination cost.
\end{corollary}

\begin{proof}
Substituting $\epsilon(n) = (T/n)^\beta$:
\[
  -\frac{T^\beta}{n^\beta} + \beta \cdot \frac{T^\beta}{n^\beta} + \ln(1-c) = 0
  \;\;\Longrightarrow\;\;
  (\beta - 1) \frac{T^\beta}{n^\beta} = -\ln(1-c) = |\ln(1-c)|.
\]
Solving for $n$:
\[
  n^\beta = \frac{(\beta-1) T^\beta}{|\ln(1-c)|}
  \;\;\Longrightarrow\;\;
  n = T \left(\frac{\beta-1}{|\ln(1-c)|}\right)^{1/\beta}.
\]
This requires $\beta > 1$ for the expression to be positive and real.
\end{proof}

%% ============================================================
\subsection{Error Amplification in Independent Pipelines}
%% ============================================================

\begin{theorem}[Error Amplification Factor]
\label{thm:error-amp}
For a uniform-capability pipeline with per-agent success probability $p$ and
coordination cost $c$, define the \emph{amplification factor} $\alpha := p(1-c)$.
Then:
\[
  P_{\mathrm{eff}}(n) = \frac{p}{1-c} \cdot \alpha^n = \frac{p \cdot \alpha^n}{1-c}.
\]
Equivalently, $P_{\mathrm{eff}}(n) = p \cdot \alpha^{n-1}$.

The system failure probability satisfies:
\[
  F_{\mathrm{sys}}(n) = 1 - p \cdot \alpha^{n-1}.
\]

For $\alpha < 1$ (which holds whenever $p < 1$ or $c > 0$ with $p \leq 1$):
\begin{enumerate}
  \item $P_{\mathrm{eff}}(n) \to 0$ exponentially as $n \to \infty$.
  \item $F_{\mathrm{sys}}(n) \to 1$ exponentially as $n \to \infty$.
  \item The \emph{reliability half-life} is $n_{1/2} = 1 + \ln 2 / |\ln \alpha|$.
\end{enumerate}
\end{theorem}

\begin{proof}
From Definition~\ref{def:coord-cost}:
\[
  P_{\mathrm{eff}}(n) = p^n (1-c)^{n-1} = p \cdot [p(1-c)]^{n-1} = p \cdot \alpha^{n-1}.
\]
Since $\alpha = p(1-c)$ and both $p \leq 1$, $(1-c) \leq 1$ with at least one strict,
we have $\alpha < 1$. Therefore $\alpha^{n-1} \to 0$ as $n \to \infty$, giving
conclusions (1) and (2).

For (3), set $P_{\mathrm{eff}}(n_{1/2}) = P_{\mathrm{eff}}(1)/2 = p/2$:
\[
  p \cdot \alpha^{n_{1/2}-1} = p/2
  \;\Longrightarrow\;
  \alpha^{n_{1/2}-1} = 1/2
  \;\Longrightarrow\;
  n_{1/2} = 1 + \frac{\ln 2}{|\ln \alpha|}.
  \qedhere
\]
\end{proof}

%% ============================================================
\subsection{Throughput-Optimal Pipeline Length}
%% ============================================================

\begin{theorem}[Universal Throughput Optimum]
\label{thm:throughput-optimum}
Define the \emph{throughput} of a uniform-capability pipeline as $\Theta(n) := n \cdot P_{\mathrm{sys}}(n) = n \, p^n$,
representing the expected number of successfully completed subtasks. Then:
\begin{enumerate}
  \item The throughput-maximizing pipeline length (treating $n$ as continuous) is
  \[
    n^* \;=\; -\frac{1}{\ln p}.
  \]
  \item At this optimum, the system success probability attains the \emph{universal constant}
  \[
    P_{\mathrm{sys}}(n^*) \;=\; e^{-1} \;\approx\; 0.3679,
  \]
  independent of $p$.
  \item The maximum throughput is $\Theta^* = n^* \cdot e^{-1} = -1/(e \ln p)$.
\end{enumerate}
\end{theorem}

\begin{proof}
For (1), differentiate $\Theta(n) = n p^n$ with respect to $n$:
\[
  \Theta'(n) = p^n + n p^n \ln p = p^n(1 + n \ln p).
\]
Setting $\Theta'(n^*) = 0$ and noting $p^{n^*} > 0$, we get $1 + n^* \ln p = 0$,
hence $n^* = -1/\ln p$. The second derivative $\Theta''(n^*) = p^{n^*} \ln p < 0$
confirms this is a maximum.

For (2), substitute $n^* = -1/\ln p$ into $P_{\mathrm{sys}} = p^n$:
\[
  P_{\mathrm{sys}}(n^*) = p^{-1/\ln p} = e^{\,\ln p \cdot (-1/\ln p)} = e^{-1}.
\]

Statement (3) follows immediately: $\Theta^* = n^* \cdot e^{-1} = -1/(e \ln p)$.
\end{proof}

\begin{remark}
The universality of $P_{\mathrm{sys}}(n^*) = e^{-1}$ means that \emph{every} throughput-optimal
pipeline, regardless of agent capability $p$, operates at the same system success probability.
Higher-capability agents simply allow longer pipelines (larger $n^*$) while maintaining
this fixed reliability. This provides a natural bridge to Regime~2: if dependent agents
can keep $P_{\mathrm{sys}} > e^{-1}$ at the same pipeline length, coordination is provably beneficial.
\end{remark}

%% ============================================================
\subsection{Correspondence with Empirical Findings}
%% ============================================================

\begin{remark}[Empirical Correspondence]
\label{rmk:empirical}
The results of this section correspond to the following empirical findings from
Kim et al.\ (2025):

\begin{enumerate}
\item \textbf{Independent coordination: $-70\%$ on sequential planning tasks.}
  Theorem~\ref{thm:critical-threshold} shows that for any $c > 0$, adding
  independent agents to a sequential pipeline always degrades performance.
  With uniform $p \approx 0.7$ and modest $c \approx 0.05$, a 3-agent pipeline gives
  $P_{\mathrm{eff}}(3) = 0.7^3 \cdot 0.95^2 \approx 0.310$, a 56\% reduction from the
  single-agent $P_{\mathrm{eff}}(1) = 0.7$.

\item \textbf{Error amplification: independent agents amplify errors $\sim\alpha^n$.}
  Theorem~\ref{thm:error-amp} derives this exactly with $\alpha = p(1-c) < 1$.

\item \textbf{Tool-coordination tradeoff: overhead compounds with environment complexity.}
  In the task decomposition model (Theorem~\ref{thm:optimal-n-general}), increasing
  task complexity $T$ shifts the optimal $n^*$ and changes the balance between
  decomposition benefit and coordination cost. The coordination overhead
  $(1-c)^{n-1}$ compounds geometrically with pipeline length.

\item \textbf{Capability saturation at $\sim 45\%$ baseline.}
  For $p = 0.45$ and $c = 0$, a 2-agent pipeline gives $0.45^2 = 0.2025$, already
  below the single-agent baseline. This prefigures the saturation result of Regime~5:
  the critical threshold $p^* = 1/(1-c)$ means sequential coordination is generically
  unprofitable, and the saturation threshold emerges from the balance with parallel
  architectures (Regime~3).
\end{enumerate}
\end{remark}

%% ============================================================
\subsection{Summary of Regime 1 Results}
%% ============================================================

We collect the principal results:

\begin{enumerate}
  \item \textbf{Multiplication Law} (Thm~\ref{thm:mult-law}):
    $P_{\mathrm{sys}} = \prod p_i$.

  \item \textbf{Exponential Decay} (Thm~\ref{thm:exp-decay}):
    $P_{\mathrm{sys}} = p^n = e^{n\ln p}$, decaying at rate $|\ln p|$ per agent.

  \item \textbf{AM--GM Bound} (Lem~\ref{lem:am-gm}):
    Heterogeneous agents are strictly worse than uniform agents with the same mean.

  \item \textbf{Critical Threshold} (Thm~\ref{thm:critical-threshold}):
    $p^* = 1/(1-c)$; for $c > 0$, $p^* > 1$, so adding agents always hurts.

  \item \textbf{Optimal Pipeline Length} (Thm~\ref{thm:optimal-n-general},
    Cor~\ref{cor:optimal-approx}):
    Under power-law decomposition with $\beta > 1$:
    $n^* = T[(\beta-1)/|\ln(1-c)|]^{1/\beta}$.

  \item \textbf{Error Amplification} (Thm~\ref{thm:error-amp}):
    $\alpha = p(1-c)$; reliability decays as $p\alpha^{n-1}$.
\end{enumerate}
